\section{Conclusion}
\label{sec:conclusion}

Three quantities are routinely conflated when a paired difference is used as a feature.
Efficiency $\eta$ compares the variance of $D$ to unpaired measurement.
Linear alignment decides whether $D$ matches $(A,B)$ in $R^2$, and under joint normality whether $D$ is sufficient for $Y$.
Reliability is decided by $\eta_S/\eta_N$, not by total $\eta$.

Fitting the pair when $D$ already aligns costs $1/(2n)$ in expected test log-loss at leading order, with a fitted second-order correction that predicts crossover sample sizes on the designs studied here.
The Gaussian $(z,r^2)$ map does not travel to heavy tails, skew, counts or an ordinal cut.
The second-moment identities do.
